\documentclass{article}
\usepackage{hyperref}

\begin{document}

\title{Software Requirements Specification (SRS)}
\author{Team 13}
\date{\today}
\maketitle

\newpage
\tableofcontents
\newpage

\section{Version Description}
Version 1.0 - Initial document setup for Snapshot 1.

\section{Introduction}

\subsection{Purpose}
The purpose of this document is to describe the software requirements for the mock Final Fantasy VII-style RPG game application. This document explains the expected functionality, user interactions, software interfaces, and legal or ethical considerations for the project.

\subsection{Intended Audience}
This document is intended for developers, testers, instructors, and team members who need to understand the project requirements and expected system behavior.

\subsection{Overview}
This project is a mock RPG game application inspired by classic turn-based role-playing games. The application will allow users to interact with a game interface, control a character, navigate menus, and participate in basic turn-based battle sequences.

\section{External Interface Requirements}

\subsection{User Interface}
The user interface should allow players to interact with the game through menus, buttons, keyboard input, or other basic controls. The interface should clearly display important game information such as player status, enemy status, available actions, and game progress.

\subsection{Software Interfaces}
The application may interact with local files for saving and loading game data. Development tools may include GitHub for version control, Jira for task tracking, TestRail for test documentation, and Docker for environment setup.

\section{Functional Requirements}

\begin{itemize}
    \item The system shall allow the user to start the game.
    \item The system shall display a main menu.
    \item The system shall allow the user to navigate the game interface.
    \item The system shall include a basic battle system.
    \item The system shall display player and enemy information during battle.
    \item The system shall allow future updates for inventory, character stats, and save data.
\end{itemize}

\section{Non-Functional Requirements}

\begin{itemize}
    \item The system should be easy to understand and use.
    \item The system should be organized in a way that allows future development.
    \item The project files should be stored and maintained using GitHub.
    \item Documentation should be updated at each project snapshot.
\end{itemize}

\section{Legal and Ethical Considerations}

\subsection{Data Privacy}
The project is not expected to collect sensitive user data. If save data is used, it should only store basic game progress such as character status, inventory, or player location.

\subsection{Ethical Issues}
The project should avoid using copyrighted assets from Final Fantasy VII directly. Any graphics, music, names, or characters should be original, placeholder, or properly credited if permitted.

\section{Glossary}

\begin{itemize}
    \item RPG - Role Playing Game
    \item UI - User Interface
    \item SRS - Software Requirements Specification
    \item GitHub - Version control and repository hosting platform
    \item Jira - Project management and task tracking tool
    \item TestRail - Test case and test report management tool
\end{itemize}

\section{References}

\begin{itemize}
    \item \href{https://docs.github.com/}{GitHub Documentation}
    \item \href{https://support.atlassian.com/jira-software-cloud/}{Jira Documentation}
    \item \href{https://support.testrail.com/}{TestRail Documentation}
    \item \href{https://www.latex-project.org/help/documentation/}{LaTeX Documentation}
\end{itemize}

\end{document}